% =====================================================================
% Appendix A.4 — Per-tap first-mode estimates
%
% Regenerated by: audit/scripts/appendix_A4_A5_tables.py  (exit 0)
% Source code reused: audit/scripts/audit_1_1_sd_vs_sem.py (tap_f1, and its
%   Table 4.5 reference values, so the cross-check is by construction)
% Audit item: 1.1
%
% CROSS-CHECK RESULT: all 13 sets reproduce Table 4.5's printed mean and tap
%   SD at printed precision. Zero disagreements. n = 5 in every cell.
%
% PACKAGES REQUIRED: booktabs
% OPTIONAL: adjustbox or a smaller font if the 9-column table overruns the
%   text block; \small is applied below and is usually sufficient at A4.
% Labels prefixed app:pertap / tab:pertap.
% =====================================================================

\section{Per-tap first-mode estimates}
\label{app:pertap}

Table~4.5 reports the mean and tap standard deviation of $f_1$ over the five taps
of each graded cell. The individual tap estimates are given here so that both
columns, and the uncertainty column derived from them, can be reproduced.

Every cell was captured as five taps on one physical damage state, so the spread
across taps measures the repeatability of the modal estimate within that state
and not the reproducibility of the damage. All thirteen sets have $n = 5$; there
are no exceptions.

\begin{table}[htbp]
\centering
\small
\caption{Per-tap first-mode estimates. Taps are in capture order. The mean and
the standard deviation as a percentage of the mean are the two columns of
Table~4.5.}
\label{tab:pertap}
\begin{tabular}{lrrrrrrrr}
\toprule
Set & Tap 1 & Tap 2 & Tap 3 & Tap 4 & Tap 5 & Mean & SD (Hz) & SD (\%) \\
\midrule
Session baseline & 2.9314 & 2.9156 & 2.9189 & 2.9208 & 2.9215 & 2.9216 & 0.0059 & 0.202 \\
\midrule
Base, trace & 2.8648 & 2.8518 & 2.8586 & 2.8534 & 2.8452 & 2.8548 & 0.0074 & 0.258 \\
Base, light & 2.1106 & 2.0462 & 2.1290 & 2.1016 & 2.0175 & 2.0810 & 0.0470 & 2.259 \\
Base, moderate & 1.3345 & 1.4459 & 1.3149 & 1.5966 & 1.7665 & 1.4917 & 0.1901 & 12.745 \\
\midrule
Floor 1, trace & 2.7971 & 2.7972 & 2.7988 & 2.7956 & 2.7950 & 2.7968 & 0.0015 & 0.054 \\
Floor 1, light & 1.7926 & 1.8915 & 1.8402 & 1.7815 & 1.7951 & 1.8202 & 0.0458 & 2.515 \\
Floor 1, moderate & 1.2752 & 1.3195 & 1.2916 & 1.2899 & 1.3074 & 1.2967 & 0.0171 & 1.318 \\
\midrule
Floor 2, trace & 2.8951 & 2.8884 & 2.8875 & 2.8972 & 2.8980 & 2.8932 & 0.0049 & 0.171 \\
Floor 2, light & 2.2293 & 2.2390 & 2.1980 & 2.2092 & 2.2084 & 2.2168 & 0.0168 & 0.759 \\
Floor 2, moderate & 1.8015 & 1.8009 & 1.8024 & 1.7988 & 1.7910 & 1.7989 & 0.0046 & 0.256 \\
\midrule
Floor 3, trace & 2.9154 & 2.9158 & 2.9236 & 2.9311 & 2.9183 & 2.9208 & 0.0066 & 0.226 \\
Floor 3, light & 2.7180 & 2.7130 & 2.7194 & 2.7194 & 2.7199 & 2.7179 & 0.0028 & 0.105 \\
Floor 3, moderate & 2.5106 & 2.5032 & 2.5081 & 2.5127 & 2.5125 & 2.5094 & 0.0039 & 0.157 \\
\bottomrule
\end{tabular}
\end{table}

\subsection*{Reproducing the uncertainty column of Table~4.5}

The uncertainty column of Table~4.5 is the standard error of the difference
between the five-tap damaged mean and the five-tap baseline mean, expressed as a
percentage of the baseline mean. From the standard deviations in
Table~\ref{tab:pertap}, with $n_b = n_t = 5$,
%
\begin{equation*}
  \mathrm{se} \;=\;
  \sqrt{\frac{s_b^{2}}{n_b} + \frac{s_t^{2}}{n_t}} \, ,
\end{equation*}
%
propagated onto $\Delta f_1 = 100\,(\bar f_t/\bar f_b - 1)$. For the base plate at
moderate damage this gives $2.91$ percentage points, the value printed in
Table~4.5. Propagating the standard deviations themselves, without dividing by
$n$, would give $6.51$ points instead; the ratio between the two is $\sqrt{5}$ in
every cell.

The base plate at moderate damage is the least repeatable cell in the campaign,
at $12.745\%$, and Table~\ref{tab:pertap} shows why: $f_1$ runs from $1.3149$ to
$1.7665$~Hz across five taps on an unchanged physical state. That behaviour is
examined in Appendix~A.1 and in Section~5.4.
